\section{Why a largest full analytic domain need not exist}\label{app:no-largest}\label{dom:section}\label{sec:domains}
In this appendix a pair $(E,F)$ consists of a source space and a value space. It is split in degree $k$
(Definition~\ref{dom:split-pair}) when $j_{E,F}$ has a bounded linear retraction, a splitting for the single
value space $F$. Split pairs form an analytic domain (Definition~\ref{dom:definition}), since a splitting at
the destination of a morphism suffices, with no compatibility between pairs (Lemma~\ref{dom:split-incoming}).

There need not be a largest analytic domain: the endomorphism actions at two objects can each be
analytic while their cross-action is not.
As in Proposition~\ref{prop:homogeneous-lift}, a bounded $d$-linear \emph{lift} of a
homogeneous map $f$ is a bounded multilinear map $B$ with $B(x,\ldots,x)=f(x)$.

\begin{definition}\label{dom:definition}
A full subcategory $\mathcal D\subseteq\mathsf{Vec}_K^{\mathrm{op}}\times\mathsf{Vec}_K$ is an
\emph{analytic domain} for $\Alt^k$ if the restriction of $\Alt^k$ to $\mathcal D$ is analytic on every
hom space.
\end{definition}

\Needspace{5\baselineskip}
\begin{definition}\label{dom:split-pair}
A pair $(E,F)$ is \emph{split in degree $k$} if the isometric inclusion
\[
j_{E,F}\colon\Alt^k(E;F)\longrightarrow\Mult^k(E;F)
\]
has a bounded linear retraction.
\end{definition}

\begin{lemma}\label{dom:split-incoming}
Every morphism-space action whose destination is a split pair is analytic. In particular, the full
subcategory of split pairs is an analytic domain.
\end{lemma}

\begin{proof}
Let $(E,F)$ be split, with retraction $R$, and let the other object be $(D,G)$. For $u\colon E\to D$
and $v\colon G\to F$, the action is
\[
m\longmapsto R\bigl(v\circ m\circ(u,\ldots,u)\bigr).
\]
Before applying $R$, replace the $k$ copies of $u$ by independent variables. This gives a bounded
$(k+1)$\nobreakdash-linear lift of the joint action, with values in
$\cL(\Alt^k(D;G),\Mult^k(E;F))$. Postcomposition with $R$ gives a lift with the required alternating
codomain. No compatibility between the retractions chosen for different objects is needed.
\end{proof}

\begin{theorem}[no largest full analytic domain]\label{dom:no-largest}
Let $p$ be prime, let $k\ge p$, and equip $K=\F_p(t)$ with its $t$\nobreakdash-adic absolute value.
There is no largest full analytic domain for
\[
\Alt^k\colon\mathsf{Vec}_K^{\mathrm{op}}\times\mathsf{Vec}_K\longrightarrow\mathsf{Vec}_K.
\]
The assertion remains true if the domains must be closed under isomorphisms, and if their objects are
restricted to pairs of finite-dimensional, nonarchimedean normed $K$\nobreakdash-spaces. In the latter case,
isomorphism closure is understood inside that restricted ambient category.

More precisely, there are objects $X_k=(E_k,G)$ and $Y_k=(D_k,G)$ whose singleton full subcategories
are analytic, such that $Y_k$ is split and the action on morphisms from $Y_k$ to $X_k$ is not analytic.
\end{theorem}

Here finite-dimensional means finite \emph{algebraic} dimension over the incomplete field $K$. These spaces
need not admit continuous coordinates over $K$. The theorem does not give a counterexample over a complete
field, or a counterexample with Banach objects.


We give the entire construction. Put $L=\F_p((t))$, the completion of $K$.
We first work in degree $p$, then extend to every $k\ge p$ by adjoining coordinate summands.

\subsection{The rigid source and the split pair}

We use the following elementary completion principle. Let $K$ be dense in a complete valued extension $L$,
let $(E_i)$ be a finite family of dense $K$\nobreakdash-subspaces of Banach $L$\nobreakdash-spaces $H_i$,
and let $Z$ be Banach over $L$.
A bounded $K$\nobreakdash-multilinear map $\prod_iE_i\to Z$ extends uniquely to a bounded
$L$\nobreakdash-multilinear map $\prod_iH_i\to Z$, with the same norm. Indeed, approximate each input by
vectors in $E_i$. A telescoping expansion and the multilinear bound make the resulting values Cauchy,
independently of the approximations. The bound persists in the limit, and density of $K$ in $L$ gives
$L$\nobreakdash-linearity in each slot. In the applications below the extended codomain is the complete
space $L$ or $A$, while membership of values on the original inputs in the smaller spaces is retained.

Put $L=\F_p((t))$, the completion of $K$, and consider the normed $L$\nobreakdash-algebra
\[
A=L[\epsilon]/(\epsilon^p),\qquad e_i=\epsilon^i\quad(0\le i<p),
\]
with the maximum norm in this basis. Multiplication satisfies $\norm{xy}\le\norm x\,\norm y$.
Choose elements $a_0,\ldots,a_{p-1}\in L$, together with the finitely many auxiliary scalars
$\tau_w$ introduced below, algebraically independent over $K$. Such a choice is possible because $K$ is
countable and $L$ is uncountable: the elements of $L$ algebraic over any finitely generated extension of
$K$ form a countable set. Set
\[
a=\sum_{i=0}^{p-1}a_i e_i,\qquad E=K^p+Ka\subseteq A.
\]
Every subspace used below carries the norm induced by $A$ or $L$. In particular, $E$ has $K$\nobreakdash-dimension
$p+1$ and completion $A$.

\begin{lemma}\label{dom:rigid}
Every bounded endomorphism of $E$ is multiplication by a scalar in $K$, and every bounded linear map
$E\to K$ is zero.
\end{lemma}

\begin{proof}
A bounded $K$\nobreakdash-linear map extends to completions, and the extension is $L$\nobreakdash-linear
by density of $K$ in $L$. If it preserves $E$, its columns, the images of the $e_i$, give a matrix
$M+ac^{\mathsf T}$, with $M\in M_p(K)$ and $c\in K^p$. Its value at $a$ lies in $K^p+Ka$.
Comparison of homogeneous degrees in the algebraically independent $a_i$ shows first that
$a(c^{\mathsf T}a)=0$, hence $c=0$, and then that $M=sI$ for some $s\in K$.

Similarly, the extension of a bounded map $E\to K$ is an $L$\nobreakdash-linear functional on $A$.
Its coefficients belong to $K$, since its values on the $e_i$ do. Its value on $a$ also lies in $K$,
so algebraic independence makes all its coefficients zero.
\end{proof}

Multiplication by $a$ has lower triangular Toeplitz matrix
\[
(M_a)_{ij}=\begin{cases}a_{i-j},&i\ge j,\\0,&i<j.\end{cases}
\]
Define
\[
D_0=\spn_K\{e_i,\epsilon^i a:0\le i<p\}+Ka^2,
\qquad
\mathcal W=\{e_i:0\le i<p\}\cup\{e_i+e_j:0\le i<j<p\},
\]
and put
\[
D=D_0+\sum_{w\in\mathcal W}K\tau_w w,\qquad
G=\spn_K\{\det(d_1,\ldots,d_p):d_i\in D\}\subseteq L.
\]
The determinant is taken in the basis $(e_0,\ldots,e_{p-1})$. These are finite-dimensional
$K$\nobreakdash-spaces, with completions $A$ and $L$: $D$ contains $K^p$, and $G$ contains $1$.
By multilinearity, $G$ is spanned by determinants of the displayed finite generating set for $D$.
We may therefore identify $G$ with a finite-dimensional subspace of
$K[a_0,\ldots,a_{p-1},(\tau_w)_w]$. Every element has degree at most one in each individual $\tau_w$,
since a nonzero determinant cannot repeat the generator $\tau_w w$. Algebraic independence gives
\begin{equation}\label{dom:tau-gap}
\tau_w^2G\cap G=\{0\}\qquad(w\in\mathcal W).
\end{equation}

\begin{lemma}\label{dom:all-alt}
Every bounded $p$\nobreakdash-linear map $D^p\to G$ is alternating. Thus $(D,G)$ is split in degree $p$.
\end{lemma}

\begin{proof}
Compose such a map $m$ with $G\hookrightarrow L$ and extend to a bounded $L$\nobreakdash-multilinear map
on $A^p$. Fix two argument positions. Put the same $w\in\mathcal W$ in those positions and arbitrary
standard basis vectors in the others, and call the value $c\in G$. Replacing both copies of $w$ by
$\tau_w w\in D$ shows that $\tau_w^2c\in G$, so~\eqref{dom:tau-gap} gives $c=0$.
Taking $w=e_i$ gives vanishing on repeated basis vectors. Taking $w=e_i+e_j$ then gives the signed
interchange identity in the two positions. Expanding two equal arbitrary arguments in the basis proves
that the completed map, and hence $m$, is alternating, including when $p=2$. The inclusion of alternating
maps into multilinear maps is therefore an isometric bijection, whose inverse is the desired retraction.
\end{proof}

\subsection{A missing coefficient}

The coefficient space of alternating maps on $E$ is
\[
C=\{\lambda\in G:a_i\lambda\in G\text{ for every }i\}.
\]
Indeed, an alternating map extends to $\lambda\det$ on the $p$\nobreakdash-dimensional completion $A$.
Determinants of tuples from $E$ span $\spn_K\{1,a_0,\ldots,a_{p-1}\}$: a determinant can contain at
most one copy of $a$, and replacing one basis vector by $a$ realizes each $a_i$ up to sign. Consequently
\begin{equation}\label{dom:coefficient-space}
\Alt^p(E;G)=C\det.
\end{equation}
This identification is isometric, since the determinant bound gives one inequality and evaluation on the
standard basis gives the reverse inequality.

Algebraically set every $\tau_w$ to zero. The image of $G$ is
\[
G_0=\spn_K\{\det(d_1,\ldots,d_p):d_i\in D_0\}.
\]
Put $C_0=\{\lambda\in G_0:a_i\lambda\in G_0\text{ for every }i\}$. Specialization sends $C$ into $C_0$;
no continuity of specialization is asserted. The displayed generators of $D_0$ are homogeneous polynomial
vectors of degrees zero, one, and two. Their determinants are homogeneous, so $G_0$ is graded, and the
defining conditions show that $C_0$ is graded as well.

\begin{lemma}\label{dom:degree-gap}
The homogeneous degree-one part of $C_0$ is zero.
\end{lemma}

\begin{proof}
The degree-two part of $G_0$ is spanned by the $2\times2$ minors of $M_a$ and the coordinates of $a^2$.
These correspond respectively to choosing two degree-one columns, or one degree-two column, in a
determinant of generators. On homogeneous quadratic polynomials define the $K$\nobreakdash-linear map
$\Phi\colon K[a_0,\ldots,a_{p-1}]_2\to K[T]$ by
\[
\Phi(a_i a_j)=
\begin{cases}
T^{i+j},&i+j\ge p-1,\\
0,&i+j<p-1.
\end{cases}
\]
Every $2\times2$ minor of $M_a$ is killed by $\Phi$. When both product terms are nonzero, their index
sums coincide. If only one is nonzero, write the rows as $r<s$ and columns as $i<j$. The surviving term
has index sum $r+s-i-j$, and the missing term implies $r<j$. Hence
$r+s-i-j\le s-i-1\le p-2$, so its image is again zero.

The coefficient of $\epsilon^r$ in $a^2$ is $\sum_{i+j=r}a_i a_j$. Its image is zero for $r<p-1$;
for $r=p-1$ there are exactly $p$ summands and its image is $pT^{p-1}=0$. Thus $\Phi$ kills the entire
degree-two part of $G_0$. If $\ell=\sum_jc_j a_j$ lies in the degree-one part of $C_0$, then
$a_{p-1}\ell\in G_0$, whereas
\[
\Phi(a_{p-1}\ell)=\sum_jc_jT^{p-1+j}.
\]
The powers are distinct, so all $c_j$ vanish.
\end{proof}

\subsection{The nonanalytic cross-action}

Set $X=(E,G)$ and $Y=(D,G)$. Lemma~\ref{dom:rigid} identifies $\cL(E,E)$ isometrically with $K$,
and the self-action of $X$ is
\[
(s,v)\longmapsto s^p v_*.
\]
This is a bounded polynomial, so the singleton full subcategory on $X$ is analytic. The singleton on $Y$
is analytic by Lemmas~\ref{dom:all-alt} and~\ref{dom:split-incoming}.

For $x\in E$, multiplication by $x$ in $A$ restricts to a bounded map $M_x\colon E\to D$.
Indeed, writing $x=b+ha$ with $b\in K^p$ and $h\in K$, its values on the generators $e_j,a$ lie in
\[
K^p+\sum_iK\epsilon^ia+Ka^2\subseteq D.
\]
Submultiplicativity gives a bounded linear family $x\mapsto M_x$ from $E$ to $\cL(E,D)$.
Its matrix is lower triangular with diagonal entry $x_0$, the $e_0$\nobreakdash-coordinate of $x$ in $A$, so $\det M_x=x_0^p$. In the cross-action from
$Y$ to $X$, fix the target morphism to $\id_G$, evaluate at $\det\in\Alt^p(D;G)$, and use
\eqref{dom:coefficient-space}. The resulting map is
\begin{equation}\label{dom:bad-coordinate}
f\colon E\longrightarrow C,\qquad f(x)=x_0^p.
\end{equation}
Its values lie in $C$ because it was obtained by pulling back an alternating $G$\nobreakdash-valued map.

Suppose that $f$ were analytic at zero. Since $f(sx)=s^pf(x)$, uniqueness of power-series coefficients
on scalar lines gives a bounded $p$\nobreakdash-linear map $B\colon E^p\to C$ with diagonal $f$.
This is the homogeneous coefficient argument used in Proposition~\ref{prop:homogeneous-lift}; it requires no factorial
division or completeness. Compose $B$ with $C\hookrightarrow L$ and extend to an $L$\nobreakdash-multilinear
map on $A^p$. Put
\[
b_i=B(e_0,\ldots,e_0,e_i)\in C.
\]
There are $p-1$ initial copies of $e_0$, and $b_0=f(e_0)=1$. Evaluation at $a$ in the last position gives
$\sum_i a_i b_i\in C$. Specializing all $\tau_w$ to zero yields elements $d_i\in C_0$ with
\[
d_0=1,\qquad \sum_i a_i d_i\in C_0.
\]
The degree-one part of this last expression is $\sum_i a_i(d_i)_0$, where $(d_i)_0\in K$ is the
constant term. Its $a_0$\nobreakdash-coefficient equals one. This contradicts Lemma~\ref{dom:degree-gap}.
Thus~\eqref{dom:bad-coordinate}, and therefore the cross-action from $Y$ to $X$, is not analytic.

\subsection{Higher degrees and the categorical conclusion}

For $k=p+r$, let
\[
E_k=E\oplus K^r,\qquad D_k=D\oplus K^r
\]
with maximum norms. Their completions have dimension $k$ over $L$. Determinants of tuples from $E_k$
still span $\spn_K\{1,a_0,\ldots,a_{p-1}\}$, so extraction of the top-degree coefficient identifies
\begin{equation}\label{dom:padded-coefficients}
\Alt^k(E_k;G)=C\det_k
\end{equation}
isometrically. In particular, the coefficient takes values in $C$, a fact needed to retain the preceding
obstruction after increasing the degree.

By Lemma~\ref{dom:rigid}, the extension of every bounded endomorphism of $E_k$ to its
$L$\nobreakdash-completion $A\oplus L^r$ has block form
\[
\begin{pmatrix}sI_p&B\\0&T\end{pmatrix},\qquad s\in K,\quad T\in M_r(K).
\]
The scalar $s$ and entries of $T$ are bounded linear functions of the endomorphism, as is seen by
evaluation on the standard basis. Its determinant is $s^p\det T$, so its action on
\eqref{dom:padded-coefficients} is multiplication by this bounded polynomial. Including postcomposition
shows that the singleton full subcategory on $X_k=(E_k,G)$ is analytic.

The pair $Y_k=(D_k,G)$ is split. Given $m\in\Mult^k(D_k;G)$, restrict the first $p$ arguments to
$D$ and fix the last $r$ arguments at the standard auxiliary basis $e'_1,\ldots,e'_r$ of $K^r$.
The resulting $p$-linear map $\alpha_m$ is alternating by Lemma~\ref{dom:all-alt}. Write
$\pi_D\colon D_k\to D$ and $\pi_r\colon D_k\to K^r$ for the two projections. Explicitly, set
\[
 (\mathcal Rm)(v_1,\ldots,v_k)=
 \sum_{\substack{S\subseteq\{1,\ldots,k\}\\|S|=p}}
 \varepsilon(S)\,
 \alpha_m(\pi_Dv_{s_1},\ldots,\pi_Dv_{s_p})
 \det(\pi_rv_{t_1},\ldots,\pi_rv_{t_r}),
\]
where $s_1<\cdots<s_p$ list $S$, $t_1<\cdots<t_r$ list its complement, and $\varepsilon(S)$
is the sign of the permutation $(s_1,\ldots,s_p,t_1,\ldots,t_r)$. The determinant and the sign
are scalars multiplying the $G$-valued factor. In degree $r=0$ the determinant is $1$.
This shuffle expression is alternating by the same repeated-argument cancellation as the wedge
product, involves no factorial division, and gives a bounded linear map
$\mathcal R\colon\Mult^k(D_k;G)\to\Alt^k(D_k;G)$ with norm at most $\binom{k}{p}$.
It fixes every alternating map: after extension to the $k$\nobreakdash-dimensional completion, both maps
are determined by their value on the standard basis, where the shuffle formula has exactly one nonzero
term. This is the required retraction. For $r>0$, we do not need or assert that every $k$\nobreakdash-linear
map on $D_k$ is alternating.

Finally, restrict the cross-action along the bounded affine family
\[
E\longrightarrow\cL(E_k,D_k),\qquad x\longmapsto M_x\oplus I_r.
\]
Evaluation at $\det_k$, which is $G$\nobreakdash-valued on $D_k$ by Laplace expansion along the first $p$ rows, and coefficient extraction through~\eqref{dom:padded-coefficients} give exactly
the nonanalytic map~\eqref{dom:bad-coordinate}. Thus the cross-action from $Y_k$ to $X_k$ is not analytic.

A largest analytic full subcategory would have to contain both individually analytic singleton
subcategories, and fullness would then require this cross-action to be analytic. This is impossible.
If closure under isomorphisms is required, replace each singleton by its full isomorphism closure.
Every action between two copies is transported from the self-action by bounded linear changes of
coordinates, so both closures remain analytic. The same contradiction proves Theorem~\ref{dom:no-largest}.

\paragraph{\textit{Finite-dimensional bases.}}
The same spaces show that Theorem~\ref{thm:finite-parameters} needs continuous coordinates on the base,
not only finite algebraic dimension. Regard $E$ as an analytic manifold with one chart; by Lemma~\ref{dom:rigid} it has no
continuous coordinates. Over it take the trivial bundle with fiber $D_k\oplus E_k$ and the two frames
$\id$ and $S_x(d,e)=\bigl(d+(M_x\oplus I_r)e,e\bigr)$, together with the trivial bundle with fiber $G$;
$S_x$ and its inverse $(d,e)\mapsto\bigl(d-(M_x\oplus I_r)e,e\bigr)$ are affine in $x$, so these bundles
are analytic. By~\eqref{eq:shear-compression}, evaluation at $\det_k$
and~\eqref{dom:padded-coefficients}, analyticity of the induced alternating atlas at $x=0$ would
make~\eqref{dom:bad-coordinate} analytic at zero, which was excluded above.
